\documentclass[11pt]{article}
\usepackage[utf8]{inputenc}
\usepackage{amsmath,amssymb,amsthm}
\usepackage{geometry}
\geometry{margin=1in}
\usepackage{hyperref}

\newtheorem{theorem}{Theorem}
\newtheorem{lemma}[theorem]{Lemma}
\newtheorem{corollary}[theorem]{Corollary}
\newtheorem{definition}[theorem]{Definition}
\newtheorem{assumption}[theorem]{Assumption}
\newtheorem{remark}[theorem]{Remark}

\title{Theoretical Analysis of Gradient-Aligned Representation Debiasing and Confidence-Based Fairness Diagnosis in Federated Learning}
\author{}
\date{}

\begin{document}
\maketitle

\section{Notation}

\begin{itemize}
    \item $\mathcal{X}$: input space; $\mathcal{Y} = \{0, 1\}$: label space; $\mathcal{G} = \{0, 1\}$: sensitive group space.
    \item $f_\phi: \mathcal{X} \to \mathbb{R}^d$: feature extractor (encoder) parameterized by $\phi$.
    \item $h_\psi: \mathbb{R}^d \to \mathbb{R}$: linear classifier head parameterized by $\psi = (w, b)$, i.e., $h_\psi(z) = w^\top z + b$.
    \item $z = f_\phi(x) \in \mathbb{R}^d$: feature representation of input $x$.
    \item $\hat{p} = \sigma(h_\psi(z))$: predicted probability, where $\sigma(t) = 1/(1+e^{-t})$.
    \item $\ell(\hat{p}, y)$: binary cross-entropy loss, $\ell(\hat{p}, y) = -[y \log \hat{p} + (1-y)\log(1-\hat{p})]$.
    \item $\mathcal{L}_{\text{task}}(x, y) = \ell(\sigma(w^\top f_\phi(x) + b), y)$: task loss for a single sample.
    \item $\nabla_z \mathcal{L}_{\text{task}} = \frac{\partial \mathcal{L}_{\text{task}}}{\partial z} \in \mathbb{R}^d$: gradient of task loss w.r.t.\ feature $z$.
\end{itemize}

\begin{definition}[Representation Bias]
    For label $l \in \{0, 1\}$, the representation bias between groups $g=0$ and $g=1$ is:
    \begin{equation}
        \Delta_{\mathrm{rep}}^{(l)} = \left\| \mathbb{E}[z \mid g=0, y=l] - \mathbb{E}[z \mid g=1, y=l] \right\|_2.
    \end{equation}
\end{definition}

\begin{definition}[Equalized Odds Gap]
    For label $l \in \{0, 1\}$, the Equalized Odds (EO) gap is:
    \begin{equation}
        \mathrm{EO}^{(l)} = \left| P(\hat{y}=1 \mid g=0, y=l) - P(\hat{y}=1 \mid g=1, y=l) \right|.
    \end{equation}
\end{definition}

\begin{definition}[Group Gradient Prototype]
    For group $g$ and label $l$, the group gradient prototype is:
    \begin{equation}
        \bar{g}_{g,l} = \mathbb{E}\left[ \nabla_z \mathcal{L}_{\text{task}}(x, y) \mid g, l \right] \in \mathbb{R}^d.
    \end{equation}
\end{definition}

\begin{definition}[Gradient Direction Consistency]
    For label $l$, the gradient direction consistency between groups is:
    \begin{equation}
        \rho_l = \frac{\bar{g}_{g0,l} \cdot \bar{g}_{g1,l}}{\|\bar{g}_{g0,l}\| \cdot \|\bar{g}_{g1,l}\|} \in [-1, 1].
    \end{equation}
\end{definition}

\begin{definition}[Confidence]
    For a sample with logit $s = w^\top z + b$, the confidence is:
    \begin{equation}
        c = \max(\hat{p}, 1-\hat{p}) = \sigma(|s|).
    \end{equation}
\end{definition}

\begin{definition}[Decision Boundary Distance]
    For a sample with feature $z$, the decision boundary distance is:
    \begin{equation}
        d = \frac{|w^\top z + b|}{\|w\|}.
    \end{equation}
\end{definition}


\section{Preliminaries}

\begin{assumption}[Linear Classifier Head]
    \label{ass:linear}
    The classifier head $h_\psi$ is linear: $h_\psi(z) = w^\top z + b$.
\end{assumption}

\begin{remark}
    This assumption is consistent with the PDFFed architecture, where prototypes are passed through the last linear layer via \texttt{only\_clf\_forward}. The encoder $f_\phi$ can be arbitrarily complex (e.g., BERT); the linearity assumption applies only to the classifier head that operates on prototypes.
\end{remark}

\begin{assumption}[Sub-Gaussian Features]
    \label{ass:subgaussian}
    For each group $g$ and label $l$, the feature distribution $z \mid g, y=l$ is sub-Gaussian with mean $\mu_{g,l}$ and covariance matrix $\Sigma_{g,l}$.
\end{assumption}

\begin{assumption}[Bounded Covariance Difference]
    \label{ass:bounded_cov}
    There exists a constant $\Delta_\Sigma$ such that:
    \begin{equation}
        \|\Sigma_{g0,l} - \Sigma_{g1,l}\|_{\mathrm{op}} \leq \Delta_\Sigma, \quad \forall l \in \{0,1\}.
    \end{equation}
    When $\Delta_\Sigma = 0$, this reduces to the shared covariance assumption.
\end{assumption}

\begin{assumption}[Bounded Feature Norm]
    \label{ass:bounded_norm}
    There exists a constant $\sigma_{\max}$ such that for all $g, l$:
    \begin{equation}
        \sqrt{w^\top \Sigma_{g,l} w} \leq \sigma_{\max}.
    \end{equation}
\end{assumption}


\section{Representation Bias and Fairness}

\begin{theorem}[Representation Bias Bounds EO Gap]
    \label{thm:rep_eo}
    Under Assumptions \ref{ass:linear}--\ref{ass:bounded_norm}, the EO gap satisfies:
    \begin{equation}
        \mathrm{EO}^{(l)} \leq \frac{\|w\|}{\sqrt{2\pi}\,\sigma_{\max}} \cdot \Delta_{\mathrm{rep}}^{(l)} + O\!\left(\frac{\|w\|^2 \Delta_\Sigma}{\sigma_{\max}^3}\right).
    \end{equation}
    When $\Delta_\Sigma = 0$ (shared covariance), this simplifies to:
    \begin{equation}
        \mathrm{EO}^{(l)} \leq \frac{\|w\|}{\sqrt{2\pi}\,\sigma} \cdot \Delta_{\mathrm{rep}}^{(l)}.
    \end{equation}
\end{theorem}

\begin{proof}
    Under the linear classifier, the predicted probability for a sample with feature $z$ is $\hat{p} = \sigma(w^\top z + b)$.

    \textbf{Step 1: Express EO in terms of feature distributions.}

    The group-conditional prediction probability is:
    \begin{equation}
        P(\hat{y}=1 \mid g, y=l) = \mathbb{E}_{z \sim \mathcal{D}_{g,l}}\left[\sigma(w^\top z + b)\right],
    \end{equation}
    where $\mathcal{D}_{g,l}$ denotes the feature distribution of group $g$, label $l$.

    \textbf{Step 2: Apply sub-Gaussian concentration.}

    Let $s = w^\top z + b$. Then $s \mid g, y=l$ is sub-Gaussian with mean $\mu'_g = w^\top \mu_{g,l} + b$ and variance $\sigma_g^2 = w^\top \Sigma_{g,l} w$.

    By the sub-Gaussian property, the CDF of $s$ is approximated by $\Phi(s/\sigma_g)$ with Berry--Esseen type error $O(1/\sqrt{n})$, where $\Phi$ is the standard normal CDF. Thus:
    \begin{equation}
        P(\hat{y}=1 \mid g, y=l) \approx \Phi\!\left(\frac{\mu'_g}{\sigma_g}\right).
    \end{equation}

    \textbf{Step 3: Bound the EO gap.}

    \begin{align}
        \mathrm{EO}^{(l)} &= \left|\Phi\!\left(\frac{\mu'_{g0}}{\sigma_{g0}}\right) - \Phi\!\left(\frac{\mu'_{g1}}{\sigma_{g1}}\right)\right|.
    \end{align}

    By the mean value theorem, there exists $\xi$ between $\mu'_{g0}/\sigma_{g0}$ and $\mu'_{g1}/\sigma_{g1}$ such that:
    \begin{equation}
        \mathrm{EO}^{(l)} = \phi(\xi) \cdot \left|\frac{\mu'_{g0}}{\sigma_{g0}} - \frac{\mu'_{g1}}{\sigma_{g1}}\right|,
    \end{equation}
    where $\phi$ is the standard normal PDF with $\phi(\xi) \leq 1/\sqrt{2\pi}$.

    \textbf{Step 4: Handle shared covariance ($\sigma_{g0} = \sigma_{g1} = \sigma$).}

    \begin{equation}
        \mathrm{EO}^{(l)} \leq \frac{1}{\sqrt{2\pi}} \cdot \frac{|\mu'_{g0} - \mu'_{g1}|}{\sigma} = \frac{1}{\sqrt{2\pi}} \cdot \frac{|w^\top(\mu_{g0,l} - \mu_{g1,l})|}{\sigma}.
    \end{equation}

    By Cauchy--Schwarz:
    \begin{equation}
        |w^\top(\mu_{g0,l} - \mu_{g1,l})| \leq \|w\| \cdot \|\mu_{g0,l} - \mu_{g1,l}\| = \|w\| \cdot \Delta_{\mathrm{rep}}^{(l)}.
    \end{equation}

    Therefore:
    \begin{equation}
        \mathrm{EO}^{(l)} \leq \frac{\|w\|}{\sqrt{2\pi}\,\sigma} \cdot \Delta_{\mathrm{rep}}^{(l)}.
    \end{equation}

    \textbf{Step 5: Handle bounded covariance difference ($\sigma_{g0} \neq \sigma_{g1}$).}

    Write $\sigma_{g0} = \sigma + \delta_0$, $\sigma_{g1} = \sigma + \delta_1$ where $|\delta_g| \leq \Delta_\Sigma/(2\sigma)$. Then:
    \begin{align}
        \left|\frac{\mu'_{g0}}{\sigma_{g0}} - \frac{\mu'_{g1}}{\sigma_{g1}}\right|
        &\leq \frac{|\mu'_{g0} - \mu'_{g1}|}{\sigma} + |\mu'_{g0}| \cdot \left|\frac{1}{\sigma_{g0}} - \frac{1}{\sigma}\right| + |\mu'_{g1}| \cdot \left|\frac{1}{\sigma_{g1}} - \frac{1}{\sigma}\right| \\
        &\leq \frac{\|w\| \cdot \Delta_{\mathrm{rep}}^{(l)}}{\sigma} + O\!\left(\frac{\|w\|^2 \Delta_\Sigma}{\sigma^3}\right).
    \end{align}

    Combining with the PDF bound:
    \begin{equation}
        \mathrm{EO}^{(l)} \leq \frac{\|w\|}{\sqrt{2\pi}\,\sigma_{\max}} \cdot \Delta_{\mathrm{rep}}^{(l)} + O\!\left(\frac{\|w\|^2 \Delta_\Sigma}{\sigma_{\max}^3}\right).
    \end{equation}
\end{proof}


\section{Gradient Alignment and Representation Bias}

\begin{lemma}[Gradient Difference Norm and Direction Consistency]
    \label{lem:grad_diff}
    Let $u_l = \bar{g}_{g0,l} - \bar{g}_{g1,l}$ be the gradient difference vector. Then:
    \begin{equation}
        \|u_l\|^2 = \|\bar{g}_{g0,l}\|^2 + \|\bar{g}_{g1,l}\|^2 - 2\|\bar{g}_{g0,l}\|\|\bar{g}_{g1,l}\|\rho_l.
    \end{equation}
    Consequently:
    \begin{enumerate}
        \item[(a)] When $\rho_l = 1$: $\|u_l\| = \big|\|\bar{g}_{g0,l}\| - \|\bar{g}_{g1,l}\|\big|$ (minimum).
        \item[(b)] When $\rho_l = -1$: $\|u_l\| = \|\bar{g}_{g0,l}\| + \|\bar{g}_{g1,l}\|$ (maximum).
    \end{enumerate}
\end{lemma}

\begin{proof}
    Direct expansion of the squared norm:
    \begin{align}
        \|u_l\|^2 &= \|\bar{g}_{g0,l} - \bar{g}_{g1,l}\|^2 \\
        &= \|\bar{g}_{g0,l}\|^2 + \|\bar{g}_{g1,l}\|^2 - 2\,\bar{g}_{g0,l} \cdot \bar{g}_{g1,l} \\
        &= \|\bar{g}_{g0,l}\|^2 + \|\bar{g}_{g1,l}\|^2 - 2\|\bar{g}_{g0,l}\|\|\bar{g}_{g1,l}\|\rho_l.
    \end{align}
    The last step uses the definition $\rho_l = (\bar{g}_{g0,l} \cdot \bar{g}_{g1,l})/(\|\bar{g}_{g0,l}\|\|\bar{g}_{g1,l}\|)$.

    For (a): $\rho_l = 1 \Rightarrow \|u_l\|^2 = (\|\bar{g}_{g0,l}\| - \|\bar{g}_{g1,l}\|)^2$.

    For (b): $\rho_l = -1 \Rightarrow \|u_l\|^2 = (\|\bar{g}_{g0,l}\| + \|\bar{g}_{g1,l}\|)^2$.
\end{proof}

\begin{theorem}[Representation Bias Dynamics Under Gradient Descent]
    \label{thm:rep_dynamics}
    Let $\mu_{g,l}^{(t)} = \mathbb{E}[z^{(t)} \mid g, y=l]$ be the group mean at step $t$, and $v_l^{(t)} = \mu_{g0,l}^{(t)} - \mu_{g1,l}^{(t)}$ be the representation bias vector. Under gradient descent with learning rate $\eta$, the squared representation bias evolves as:
    \begin{equation}
        \left(\Delta_{\mathrm{rep}}^{(l,t+1)}\right)^2 - \left(\Delta_{\mathrm{rep}}^{(l,t)}\right)^2 = -2\eta\, v_l^{(t)} \cdot u_l + \eta^2 \|u_l\|^2,
    \end{equation}
    where $u_l = \bar{g}_{g0,l} - \bar{g}_{g1,l}$.

    Furthermore, by Lemma~\ref{lem:grad_diff}, the second term satisfies:
    \begin{equation}
        \eta^2 \|u_l\|^2 = \eta^2\!\left(\|\bar{g}_{g0,l}\|^2 + \|\bar{g}_{g1,l}\|^2 - 2\|\bar{g}_{g0,l}\|\|\bar{g}_{g1,l}\|\rho_l\right).
    \end{equation}
    This term is monotonically decreasing in $\rho_l$.
\end{theorem}

\begin{proof}
    \textbf{Step 1: Feature update rule.}

    Under gradient descent, for a sample $(x, y, g)$:
    \begin{equation}
        z^{(t+1)} = z^{(t)} - \eta \cdot \nabla_z \mathcal{L}_{\text{task}}(x, y).
    \end{equation}

    \textbf{Step 2: Group mean update.}

    Taking expectation over group $g$, label $l$:
    \begin{equation}
        \mu_{g,l}^{(t+1)} = \mu_{g,l}^{(t)} - \eta \cdot \bar{g}_{g,l}.
    \end{equation}

    \textbf{Step 3: Representation bias vector update.}

    \begin{align}
        v_l^{(t+1)} &= \mu_{g0,l}^{(t+1)} - \mu_{g1,l}^{(t+1)} \\
        &= (\mu_{g0,l}^{(t)} - \eta \bar{g}_{g0,l}) - (\mu_{g1,l}^{(t)} - \eta \bar{g}_{g1,l}) \\
        &= v_l^{(t)} - \eta u_l.
    \end{align}

    \textbf{Step 4: Squared norm expansion (exact, no inequality).}

    \begin{align}
        \|v_l^{(t+1)}\|^2 &= \|v_l^{(t)} - \eta u_l\|^2 \\
        &= \|v_l^{(t)}\|^2 - 2\eta\, v_l^{(t)} \cdot u_l + \eta^2 \|u_l\|^2.
    \end{align}

    Therefore:
    \begin{equation}
        \left(\Delta_{\mathrm{rep}}^{(l,t+1)}\right)^2 - \left(\Delta_{\mathrm{rep}}^{(l,t)}\right)^2 = -2\eta\, v_l^{(t)} \cdot u_l + \eta^2 \|u_l\|^2.
    \end{equation}

    \textbf{Step 5: Substitute $\|u_l\|^2$ from Lemma~\ref{lem:grad_diff}.}

    \begin{equation}
        \eta^2 \|u_l\|^2 = \eta^2\!\left(\|\bar{g}_{g0,l}\|^2 + \|\bar{g}_{g1,l}\|^2 - 2\|\bar{g}_{g0,l}\|\|\bar{g}_{g1,l}\|\rho_l\right).
    \end{equation}

    The derivative with respect to $\rho_l$ is:
    \begin{equation}
        \frac{\partial (\eta^2 \|u_l\|^2)}{\partial \rho_l} = -2\eta^2 \|\bar{g}_{g0,l}\|\|\bar{g}_{g1,l}\| < 0.
    \end{equation}

    Therefore $\eta^2\|u_l\|^2$ is strictly decreasing in $\rho_l$.
\end{proof}

\begin{remark}
    Theorem~\ref{thm:rep_dynamics} provides an \emph{exact} expression (not a bound) for the change in squared representation bias. The change depends on two terms:
    \begin{itemize}
        \item \textbf{Directional term} $(-2\eta\, v \cdot u)$: depends on the alignment between the current bias direction $v$ and the gradient difference $u$. This term can be positive or negative.
        \item \textbf{Magnitude term} $(\eta^2 \|u\|^2)$: always non-negative, and is monotonically decreasing in $\rho_l$. This term represents the ``inertia'' of bias change due to gradient magnitude.
    \end{itemize}
    The gradient alignment regularization $\mathcal{L}_{\mathrm{grad}} = 1 - \rho_l$ directly minimizes the magnitude term by maximizing $\rho_l$, thereby constraining the rate at which representation bias can change.
\end{remark}

\begin{lemma}[Convergence of Gradient Alignment Regularization]
    \label{lem:convergence}
    Let the total loss be $\mathcal{L} = \mathcal{L}_{\text{task}} + \lambda \mathcal{L}_{\mathrm{grad}}$, where $\mathcal{L}_{\mathrm{grad}} = 1 - \rho_l$ and $\lambda > 0$. If:
    \begin{enumerate}
        \item $\mathcal{L}_{\text{task}}$ is $L$-smooth,
        \item $\mathcal{L}_{\mathrm{grad}}$ is $M$-smooth,
        \end{enumerate}
    then for learning rate $\eta \leq 1/(L + \lambda M)$, gradient descent on $\mathcal{L}$ converges:
    \begin{equation}
        \mathcal{L}^{(t+1)} - \mathcal{L}^{(t)} \leq -\eta\!\left(1 - \frac{\eta(L+\lambda M)}{2}\right) \|\nabla \mathcal{L}^{(t)}\|^2.
    \end{equation}
\end{lemma}

\begin{proof}
    This is a standard result in smooth optimization. Since $\mathcal{L}_{\text{task}}$ is $L$-smooth and $\mathcal{L}_{\mathrm{grad}}$ is $M$-smooth, the total loss $\mathcal{L}$ is $(L + \lambda M)$-smooth. By the descent lemma:
    \begin{equation}
        \mathcal{L}^{(t+1)} \leq \mathcal{L}^{(t)} + \nabla \mathcal{L}^{(t)} \cdot (-\eta \nabla \mathcal{L}^{(t)}) + \frac{(L+\lambda M)\eta^2}{2} \|\nabla \mathcal{L}^{(t)}\|^2.
    \end{equation}
    Rearranging:
    \begin{equation}
        \mathcal{L}^{(t+1)} - \mathcal{L}^{(t)} \leq -\eta\!\left(1 - \frac{\eta(L+\lambda M)}{2}\right) \|\nabla \mathcal{L}^{(t)}\|^2.
    \end{equation}
    When $\eta \leq 1/(L+\lambda M)$, the coefficient is non-negative, ensuring descent.
\end{proof}

\begin{theorem}[Main Theorem: Gradient Alignment Constrains EO Gap]
    \label{thm:main}
    Under Assumptions \ref{ass:linear}--\ref{ass:bounded_norm}, with gradient alignment regularization $\mathcal{L}_{\mathrm{grad}} = 1 - \rho_l$ and learning rate satisfying Lemma~\ref{lem:convergence}, the EO gap at convergence satisfies:
    \begin{equation}
        \mathrm{EO}^{(l, \infty)} \leq \frac{\|w\|}{\sqrt{2\pi}\,\sigma_{\max}} \cdot \Delta_{\mathrm{rep}}^{(l, \infty)} + O\!\left(\frac{\|w\|^2 \Delta_\Sigma}{\sigma_{\max}^3}\right),
    \end{equation}
    where $\Delta_{\mathrm{rep}}^{(l, \infty)}$ is the representation bias at convergence. The gradient alignment regularization constrains $\Delta_{\mathrm{rep}}^{(l, \infty)}$ by:
    \begin{enumerate}
        \item Maximizing $\rho_l$, which minimizes $\|u_l\|$ (Lemma~\ref{lem:grad_diff}), thereby minimizing the magnitude term $\eta^2\|u_l\|^2$ in the bias dynamics (Theorem~\ref{thm:rep_dynamics}).
        \item Ensuring convergence of the total loss (Lemma~\ref{lem:convergence}), so the system reaches a stable state where the bias dynamics are bounded.
    \end{enumerate}
\end{theorem}

\begin{proof}
    By Theorem~\ref{thm:rep_eo}, at any step $t$:
    \begin{equation}
        \mathrm{EO}^{(l,t)} \leq C \cdot \Delta_{\mathrm{rep}}^{(l,t)} + \epsilon,
    \end{equation}
    where $C = \|w\|/(\sqrt{2\pi}\,\sigma_{\max})$ and $\epsilon = O(\|w\|^2 \Delta_\Sigma / \sigma_{\max}^3)$.

    By Lemma~\ref{lem:convergence}, the total loss $\mathcal{L} = \mathcal{L}_{\text{task}} + \lambda \mathcal{L}_{\mathrm{grad}}$ converges to a stationary point where $\nabla \mathcal{L} \approx 0$.

    At convergence, the gradient $\bar{g}_{g,l} \to 0$ for all $g, l$, so $u_l \to 0$ and the bias dynamics (Theorem~\ref{thm:rep_dynamics}) stabilize: $\left(\Delta_{\mathrm{rep}}^{(l,t+1)}\right)^2 - \left(\Delta_{\mathrm{rep}}^{(l,t)}\right)^2 \to 0$.

    During training, the gradient alignment regularization ensures that $\rho_l$ is maximized, which minimizes $\|u_l\|$ at each step (Lemma~\ref{lem:grad_diff}). A smaller $\|u_l\|$ means:
    \begin{itemize}
        \item The magnitude term $\eta^2\|u_l\|^2$ is smaller, limiting the growth of $\Delta_{\mathrm{rep}}$.
        \item The directional term $-2\eta\, v \cdot u$ is also smaller in magnitude (since $\|u_l\|$ is smaller), limiting both growth and reduction of $\Delta_{\mathrm{rep}}$.
    \end{itemize}
    Therefore, the representation bias $\Delta_{\mathrm{rep}}^{(l,\infty)}$ at convergence is constrained by the gradient alignment, and by Theorem~\ref{thm:rep_eo}, the EO gap is correspondingly bounded.
\end{proof}


\section{Confidence-Based Fairness Diagnosis}

\begin{theorem}[Confidence and Decision Boundary Distance]
    \label{thm:conf_dist}
    Under Assumption~\ref{ass:linear}, the confidence $c$ and decision boundary distance $d$ satisfy:
    \begin{equation}
        c = \sigma(\|w\| \cdot d).
    \end{equation}
    Therefore, $c$ is strictly monotonically increasing in $d$.
\end{theorem}

\begin{proof}
    By definition:
    \begin{equation}
        c = \sigma(|s|) = \sigma(|w^\top z + b|), \quad d = \frac{|w^\top z + b|}{\|w\|}.
    \end{equation}
    Therefore $|w^\top z + b| = \|w\| \cdot d$, and:
    \begin{equation}
        c = \sigma(\|w\| \cdot d).
    \end{equation}
    Since $\sigma$ is strictly increasing, $c$ is strictly increasing in $d$.
\end{proof}

\begin{lemma}[Decision Boundary Distance and Group Prediction Probability]
    \label{lem:dist_prob}
    Under Assumptions \ref{ass:linear}--\ref{ass:subgaussian}, for group $g$ and label $l$:
    \begin{equation}
        \mathbb{E}[d \mid g, y=l] = \frac{|\mu'_g|}{\|w\|} + O(\sigma_g / \|w\|),
    \end{equation}
    where $\mu'_g = w^\top \mu_{g,l} + b$ and $\sigma_g = \sqrt{w^\top \Sigma_{g,l} w}$.

    Furthermore, the group prediction probability satisfies:
    \begin{equation}
        P(\hat{y}=1 \mid g, y=l) = \Phi\!\left(\frac{\mu'_g}{\sigma_g}\right) + O(1/\sqrt{n}).
    \end{equation}
\end{lemma}

\begin{proof}
    Since $s = w^\top z + b$ is sub-Gaussian with mean $\mu'_g$ and variance $\sigma_g^2$:
    \begin{equation}
        d = \frac{|s|}{\|w\|}, \quad \mathbb{E}[d] = \frac{\mathbb{E}[|s|]}{\|w\|}.
    \end{equation}
    For sub-Gaussian $s$ with mean $\mu'_g$, $\mathbb{E}[|s|] \approx |\mu'_g| + O(\sigma_g)$, giving the first result.

    For the prediction probability, by sub-Gaussian concentration (Berry--Esseen type approximation):
    \begin{equation}
        P(\hat{y}=1 \mid g, y=l) = P(s > 0) = 1 - \Phi\!\left(\frac{-\mu'_g}{\sigma_g}\right) = \Phi\!\left(\frac{\mu'_g}{\sigma_g}\right) + O(1/\sqrt{n}).
    \end{equation}
\end{proof}

\begin{theorem}[Confidence Gap and EO Gap]
    \label{thm:conf_eo}
    Under Assumptions \ref{ass:linear}--\ref{ass:bounded_cov}, define the group confidence gap:
    \begin{equation}
        \Delta_c^{(l)} = \left|\mathbb{E}[c \mid g=0, y=l] - \mathbb{E}[c \mid g=1, y=l]\right|.
    \end{equation}
    Then:
    \begin{equation}
        \Delta_c^{(l)} \leq \frac{\|w\|}{4} \cdot \left|\mathbb{E}[d \mid g=0, y=l] - \mathbb{E}[d \mid g=1, y=l]\right|,
    \end{equation}
    and the EO gap satisfies:
    \begin{equation}
        \mathrm{EO}^{(l)} \geq \left|\Phi\!\left(\frac{\mu'_{g0}}{\sigma_{g0}}\right) - \Phi\!\left(\frac{\mu'_{g1}}{\sigma_{g1}}\right)\right|.
    \end{equation}
    When $\mu'_{g0} \approx 0$ (low confidence group) and $|\mu'_{g1}| \gg 0$ (high confidence group), the EO gap is lower bounded by a quantity proportional to the confidence gap.
\end{theorem}

\begin{proof}
    \textbf{Step 1: Bound $\Delta_c$ using Lipschitz continuity.}

    By Theorem~\ref{thm:conf_dist}, $c = \sigma(\|w\| \cdot d)$. The sigmoid function $\sigma$ is $1/4$-Lipschitz. Therefore:
    \begin{align}
        \Delta_c^{(l)} &= |\sigma(\|w\| \cdot \mathbb{E}[d \mid g=0]) - \sigma(\|w\| \cdot \mathbb{E}[d \mid g=1])| \\
        &\leq \frac{1}{4} \cdot \|w\| \cdot |\mathbb{E}[d \mid g=0] - \mathbb{E}[d \mid g=1]|.
    \end{align}
    (This uses Jensen's inequality to move the expectation inside $\sigma$.)

    \textbf{Step 2: Relate decision boundary distance to prediction probability.}

    By Lemma~\ref{lem:dist_prob}:
    \begin{equation}
        P(\hat{y}=1 \mid g, y=l) = \Phi\!\left(\frac{\mu'_g}{\sigma_g}\right).
    \end{equation}
    The EO gap is:
    \begin{equation}
        \mathrm{EO}^{(l)} = \left|\Phi\!\left(\frac{\mu'_{g0}}{\sigma_{g0}}\right) - \Phi\!\left(\frac{\mu'_{g1}}{\sigma_{g1}}\right)\right|.
    \end{equation}

    \textbf{Step 3: Analyze the case of systematic confidence difference.}

    Suppose group $g=0$ has systematically low confidence: $\mathbb{E}[c \mid g=0] \leq c_{\text{low}}$, meaning $\mu'_{g0} \approx 0$ (samples near decision boundary).

    Suppose group $g=1$ has systematically high confidence: $\mathbb{E}[c \mid g=1] \geq c_{\text{high}}$, meaning $|\mu'_{g1}| \gg 0$ (samples far from decision boundary).

    Then:
    \begin{equation}
        P(\hat{y}=1 \mid g=0, y=l) \approx \Phi(0) = 0.5,
    \end{equation}
    \begin{equation}
        P(\hat{y}=1 \mid g=1, y=l) \approx \Phi\!\left(\frac{\mu'_{g1}}{\sigma_{g1}}\right) \approx 0 \text{ or } 1.
    \end{equation}

    Therefore:
    \begin{equation}
        \mathrm{EO}^{(l)} \approx |0.5 - 0| = 0.5 \text{ or } |0.5 - 1| = 0.5.
    \end{equation}

    More precisely, using the confidence gap $\Delta_c$ and the Lipschitz bound from Step 1:
    \begin{equation}
        \Delta_d \geq \frac{4 \Delta_c}{\|w\|},
    \end{equation}
    where $\Delta_d = |\mathbb{E}[d \mid g=0] - \mathbb{E}[d \mid g=1]|$.

    By Lemma~\ref{lem:dist_prob}, $\Delta_d \approx |\mu'_{g0} - \mu'_{g1}|/\|w\|$, and the EO gap is proportional to $|\mu'_{g0}/\sigma_{g0} - \mu'_{g1}/\sigma_{g1}|$.

    When $\mu'_{g0} \approx 0$:
    \begin{equation}
        \mathrm{EO}^{(l)} \approx \Phi\!\left(\frac{|\mu'_{g1}|}{\sigma_{g1}}\right) - 0.5.
    \end{equation}
    Since $|\mu'_{g1}| = \|w\| \cdot \mathbb{E}[d \mid g=1]$ and $\mathbb{E}[d \mid g=1] \geq \Delta_d/2$:
    \begin{equation}
        \mathrm{EO}^{(l)} \gtrsim \Phi\!\left(\frac{\|w\| \Delta_d}{2\sigma_{g1}}\right) - 0.5 \geq \Phi\!\left(\frac{2\Delta_c}{\sigma_{g1}}\right) - 0.5.
    \end{equation}
    This shows that a larger confidence gap $\Delta_c$ leads to a larger lower bound on the EO gap.
\end{proof}

\begin{remark}
    Theorem~\ref{thm:conf_eo} establishes a \emph{qualitative} link: systematic confidence differences between groups are indicative of EO gaps. The quantitative bound depends on distributional parameters ($\sigma_g$, $\mu'_g$) that vary by dataset and model. The key insight is:
    \begin{itemize}
        \item Low confidence $\Leftrightarrow$ near decision boundary $\Leftrightarrow$ prediction probability near $0.5$.
        \item High confidence $\Leftrightarrow$ far from decision boundary $\Leftrightarrow$ prediction probability near $0$ or $1$.
        \item Systematic between-group confidence difference $\Rightarrow$ systematic prediction probability difference $\Rightarrow$ EO gap.
    \end{itemize}
\end{remark}


\section{Summary of Proof Chain}

\subsection*{Part I: Gradient Alignment $\to$ Representation Bias $\to$ Fairness}

\begin{enumerate}
    \item \textbf{Theorem~\ref{thm:rep_eo}}: $\mathrm{EO} \leq C \cdot \Delta_{\mathrm{rep}} + \epsilon$ \\
    (Representation bias bounds EO gap, under linear head + sub-Gaussian assumptions.)

    \item \textbf{Lemma~\ref{lem:grad_diff}}: $\|u_l\|^2 = \|g_0\|^2 + \|g_1\|^2 - 2\|g_0\|\|g_1\|\rho_l$ \\
    (Algebraic identity linking gradient difference norm to direction consistency.)

    \item \textbf{Theorem~\ref{thm:rep_dynamics}}: $\Delta^{(t+1)2} - \Delta^{(t)2} = -2\eta\, v \cdot u + \eta^2\|u\|^2$ \\
    (Exact dynamics of representation bias; magnitude term is decreasing in $\rho_l$.)

    \item \textbf{Lemma~\ref{lem:convergence}}: Total loss with gradient alignment regularization converges. \\
    (Standard smooth optimization result.)

    \item \textbf{Theorem~\ref{thm:main} (Main)}: Gradient alignment constrains $\Delta_{\mathrm{rep}}$, which bounds EO. \\
    (Combines Theorems \ref{thm:rep_eo}, \ref{thm:rep_dynamics}, and Lemma \ref{lem:convergence}.)
\end{enumerate}

\subsection*{Part II: Confidence $\to$ Decision Boundary $\to$ Fairness}

\begin{enumerate}
    \item \textbf{Theorem~\ref{thm:conf_dist}}: $c = \sigma(\|w\| \cdot d)$ \\
    (Confidence is monotonically increasing in decision boundary distance.)

    \item \textbf{Lemma~\ref{lem:dist_prob}}: $\mathbb{E}[d] \approx |\mu'_g|/\|w\|$, $P(\hat{y}=1) \approx \Phi(\mu'_g/\sigma_g)$ \\
    (Links decision boundary distance to prediction probability.)

    \item \textbf{Theorem~\ref{thm:conf_eo}}: $\Delta_c \leq \frac{\|w\|}{4} \Delta_d$, and large $\Delta_c \Rightarrow$ large EO lower bound. \\
    (Combines Theorem \ref{thm:conf_dist} and Lemma \ref{lem:dist_prob}.)
\end{enumerate}

\end{document}